\documentclass[12pt]{article}
\usepackage[margin=0.8in]{geometry}
\usepackage[utf8]{inputenc}
\usepackage{graphicx}
\usepackage[export]{adjustbox}
\usepackage{amsmath}
\usepackage{amsfonts}
\usepackage{amssymb}
\usepackage[version=4]{mhchem}
\usepackage{stmaryrd}
\usepackage{bbold}
%\usepackage[fallback]{xeCJK}
\usepackage{polyglossia}
\usepackage{fontspec}
\usepackage{censor}
%\setCJKmainfont{Noto Serif CJK TC}
\graphicspath{ {../Images/} }
\usepackage{enumitem}
\usepackage{tikz}
\usepackage{pgfplots}
\pgfplotsset{compat=1.18}
\usepackage{amsthm}
\usepackage{tcolorbox}
\usepackage{array}
\usepackage{tabularray}
\usepackage{makecell}
\usepackage{esvect}
%\usepackage{changepage}   % for the adjustwidth environment
\usepackage{enumitem}
\usepackage{tasks}

\tcbuselibrary{theorems}

\definecolor{GrayCen}{HTML}{d2d2d2}
\definecolor{GrayBG}{HTML}{d9d9d9}
\definecolor{BlueDef}{HTML}{104e8b}
\definecolor{GreenProp}{HTML}{025540}
\definecolor{RedMet}{HTML}{cd5556}
\definecolor{BlackSav}{HTML}{000000}
\def\censorcolor{GrayCen}
\let\svcensorrule\censorrule
\renewcommand\censorrule[1]{%
\textcolor{\censorcolor}{\svcensorrule{#1}}}

\newtcbtheorem
[]% init options
{definition}% name
{Définition}% title
{%
  colback=GrayBG,
  colframe=BlueDef,
  fonttitle=\bfseries,
}% options
{def}% prefix

\newtcbtheorem
[]% init options
{proposition}% name
{Proposition}% title
{%
  colback=GrayBG,
  colframe=GreenProp,
  fonttitle=\bfseries,
}% options
{prop}% prefix

\newtcbtheorem
[]% init options
{method}% name
{Méthode}% title
{%
  theorem name,
  colback=GrayBG,
  colframe=RedMet,
  fonttitle=\bfseries,
}% options
{met}% prefix

\newtcbtheorem
[]% init options
{savoir}% name
{Savoir-faire du chapitre}% title
{%
  theorem name,
  colback=GrayBG,
  colframe=BlackSav,
  fonttitle=\bfseries,
}% options
{sav}% prefix

\newtheoremstyle{break}
{\topsep}{\topsep}%
{}{}%
{\bfseries}{}%
{\newline}{}%
\theoremstyle{break}
\newtheorem*{example}{Exemple.}
\newtheorem*{examples}{Exemples.}
\newtheorem*{remark}{Remarque.}

\newenvironment{demo}[1][\proofname]{%
  \begin{proof}[#1]$ $\newline
  }{%
  \end{proof}
}

\setmainlanguage{french}
%\setmainfont{CMU Serif}

\title{
  \vspace{4em}
  \centering
  \tcbset{colframe=BlackSav,colback=GrayBG}
  \tcbox[tikznode]{
    Chapitre 11 \\Variables aléatoires
  }
}

\author{}
\date{}

\setlength\parindent{0pt}
%\setlist[itemize]{topsep=0pt}
\setlist{nolistsep}
%\setlist[itemize]{noitemsep}
%\setlist[enumerate]{noitemsep}

\newcolumntype{C}[1]{>{\centering\arraybackslash}m{#1}}
\newcolumntype{L}[1]{>{\raggedright\arraybackslash}m{#1}}
\newcolumntype{N}{@{}m{0pt}@{}}

\begin{document}
\maketitle

\tableofcontents

\newcommand{\hideM}[1]{\mbox{\vphantom{$#1$}\smash{\colorbox{gray!50}{$\displaystyle #1$}}} }
\newcommand{\hide}[1]{\mbox{\vphantom{#1}\smash{\colorbox{gray!50}{#1}}}}
% \newcommand{\hideM}[1]{\mbox{\vphantom{$#1$}\smash{\colorbox{gray!50}{\phantom{$\displaystyle #1$}}}} }
% \newcommand{\hide}[1]{\mbox{\vphantom{#1}\smash{\colorbox{gray!50}{\phantom{#1}}}}}

\newpage

\section{Rappels de Seconde (Indicateurs statistiques)}
Dans l'ensemble de cette partie, on considère l'exemple de la série statistiques suivante, donnant les distances des planètes du système solaire au Soleil :

\begin{center}
  \begin{tabular}{|l|l|}
    \hline
    Planète & Distance au soleil en millions de km \\
    \hline
    Mercure & 58 \\
    \hline
    Vénus & 108 \\
    \hline
    Terre & 150 \\
    \hline
    Mars & 228 \\
    \hline
    Jupiter & 778 \\
    \hline
    Saturne & 1429 \\
    \hline
    Uranus & 2871 \\
    \hline
    Neptune & 4503 \\
    \hline
  \end{tabular}
\end{center}

\subsection{Indicateurs de position}
\subsection{Moyenne}
\begin{definition}{}{}
  La moyenne d'une série statistique $x_{1}, \ldots, x_{n}$ se note $\bar{x}$ et est donné par la formule suivante :

  $$
  \bar{x}=\frac{x_{1}+\ldots+x_{n}}{n}=\hideM{\frac{\sum_{i=1}^{n} x_{i}}{n}}
  $$
\end{definition}

\begin{example}
  En moyenne, la distance d'une planète du système solaire au Soleil est d'environ 1266 millions de km.

  $$
  \bar{x}=\frac{58+108+150+228+778+1429+2871+4503}{8} \simeq 1266
  $$
\end{example}

\subsection{Médiane}
\begin{definition}{}{}
  La médiane d'une série statistique $x_{1}, \ldots, x_{n}$ est une valeur qui permet de couper l'ensemble des valeurs en deux parties égales : la moitié des valeurs étant plus grande que la médiane et l'autre moitié étant moins grande.
\end{definition}

\begin{remark}
  \begin{itemize}
    \item[]
    \item Lorsque la série comprend un nombre impair de valeurs, la médiane est simplement une des valeurs de la série.
    \item Lorsque la série comprend un nombre pair de valeurs, la médiane est la moyenne entre les deux valeurs du «milieu».
  \end{itemize}
\end{remark}

\begin{example}
  Dans le cas des distances planète-Soleil, la médiane est la moyenne entre la $4^{e}$ et la $5^{e}$ plus grande valeur étant donné que la série comporte 8 valeurs.

  $$
  m=\frac{228+778}{2}=503 \text { millions de } \mathrm{km}
  $$
\end{example}

\subsection{Quartiles}
\begin{definition}{}{}
  \begin{itemize}
    \item Le \textbf{premier quartile} d'une série statistique est une valeur, notée $Q_{1}$ telle qu'au moins \hide{$25 \%$} des valeurs soient plus petites que $Q_{1}$.
    \item Le  \textbf{troisième quartile} d'une série statistique est une valeur, notée $Q_{3}$ telle qu'au moins \hide{$75 \%$} des valeurs soient plus petites que $Q_{3}$.
  \end{itemize}
\end{definition}

\begin{example}
  Dans le cas des distances planète-Soleil, $Q_{1}$ correspond à la moyenne entre la deuxième et la troisième valeur (car $8 \times \frac{1}{4}=2$ ) et $Q_{3}$ correspond à lamoyenne entre la sixième et la sixième valeur (car $8 \times \frac{3}{4}=6$ ).

  $$
  \begin{gathered}
    Q_{1}=129 \\
    Q_{3}=2150
  \end{gathered}
  $$
\end{example}

\begin{remark}
  L'intervalle interquartile est l'intervale $\left[Q_{1} ; Q_{3}\right]$. Il contient donc $50 \%$ des valeurs de la série statistique.
\end{remark}

\subsection{Valeur minimale et maximale}
\begin{definition}{}{}
  La \textbf{valeur minimale} d'une série statistique est la plus petite valeur de la série. la \textbf{valeur maximale} est la plus grande valeur de la série.
\end{definition}

\begin{example}
  Dans le cas des distances planète-Soleil, la valeur minimale est 58 et la valeur maximale est 4503.
\end{example}

\subsection{Indicateurs de dispersion}
\subsection{Étendue}
\begin{definition}{}{}
  L'\textbf{étendue} d'une série statistique est la différence entre la valeur maximale et la valeur minimale.
\end{definition}

\begin{example}
  Dans le cas des distances planète-Soleil, l'étendue est de $4503-58=4445$ millions de km.
\end{example}

\subsection{Écart interquartile}
\begin{definition}{}{}
  L'\textbf{écart interquartile} est \hide{$Q_{3}-Q_{1}$}.
\end{definition}

\begin{example}
  Dans le cas des distances planète-Soleil, l'écart interquartile est $2150-129=2021$ millions de km .
\end{example}

\subsection{Écart-type}
\begin{definition}{}{}
  \begin{itemize}
    \item La \textbf{variance} $V$ d'une série statistique $x_{1}, \ldots, x_{n}$ est définie par :

      $$
      V=\frac{\left(x_{1}-\bar{x}\right)^{2}+\ldots+\left(x_{n}-\bar{x}\right)^{2}}{n}=\hideM{\frac{\sum_{i=1}^{n}\left(x_{i}-\bar{x}\right)^{2}}{n}}
      $$

    \item L'\textbf{écart type} $\sigma$ d'une série statistique est définie par :

      $$
      \sigma=\hideM{\sqrt{V}}
      $$
  \end{itemize}
\end{definition}

\begin{remark}
  On retiendra que la variance est «la moyenne des carrés des écarts à la moyenne »
\end{remark}

\begin{example}
  Dans le cas des distances planète-Soleil, on a :\\
  $\bar{x}=1266$\\
  $V=\dfrac{(58-1266)^{2}+(108-1266)^{2}+(150-1266)^{2}+\ldots+(4503-1266)^{2}}{8} \simeq 2304324$\\
  $\sigma=\sqrt{V} \simeq 1518$
\end{example}

\begin{remark}
  La donnée de la médiane et de l'écart interquartile n'est pas sensible aux valeurs extrêmes, contrairement à la donnée de la moyenne et de l'écart-type.
\end{remark}

\section{Échantillonnage}
\begin{proposition}{(admise)}{}
  Si l'effectif d'une série statistique est suffisamment grand, et si le phénomène étudié suit une loi de Gauss (la répartition des valeurs se fait selon une courbe en cloche) :

  \begin{itemize}
    \item L'intervalle $[\bar{x}-\sigma ; \bar{x}+\sigma]$ contient environ $68 \%$ des valeurs.
    \item L'intervalle $[\bar{x}-2 \sigma ; \bar{x}+2 \sigma]$ contient environ $95 \%$ des valeurs.
    \item L'intervalle $[\bar{x}-3 \sigma ; \bar{x}+3 \sigma]$ contient environ $99 \%$ des valeurs.
  \end{itemize}
\end{proposition}

\newpage
\begin{examples}
  Les phénomènes suivants suivent une loi de Gauss :

  \begin{itemize}
    \item Pourcentage de Pile obtenus lors du lancer de 50 pièces;
    \item La taille des femmes dans la population française;
    \item Le nombre de votants par bureau de vote lors d'une élection;
    \item La taille des tiges des coquelicots dans un champ;
    \item Les notes d'une évaluation.
  \end{itemize}

  Ainsi, si la moyenne a une évaluation est $\bar{x}=10$ et l'écart-type est $\sigma=3$, cela signifie qu'environ $68 \%$ des notes sont comprises entre 7 et 13 .
\end{examples}

\begin{proposition}{Loi des Grands Nombres (admise)}{}
On répète $n$ fois une même expérience aléatoire de façon indépendante et dont la probabilité de succès est $p$. Expérimentalement, plus $n$ sera grand, plus la fréquence du succès se rapprochera de $p$. De plus, on peut considérer que la fréquence de succès suit une loi de Gauss et on a :

$$
\sigma \simeq \hideM{\sqrt{\frac{p(1-p)}{n}}}
$$
  
\end{proposition}

\begin{remark}
  En général, lorsque $n \geqslant 25$, on considère que $n$ est suffisament grand pour utiliser la loi des Grands Nombres.
\end{remark}

\begin{example}
  On lance 200 fois un dé équilibré à six faces et on compte le nombre de 6 obtenus. Déterminer avec une marge d'erreur de $5 \%$, un intervalle contenant la fréquence de 6 obtenus.

  ~\\Solution :\\
  Les lancers étant indépendants les uns des autres, on peut utiliser la loi des Grands Nombres. On a $p=\frac{1}{6}$ donc :

  $$
  \sigma \simeq \sqrt{\frac{\frac{1}{6}\left(1-\frac{1}{6}\right)}{200}} \simeq 0,0264
  $$

  Par conséquent, cela signifie que l'intervalle suivant contient $95 \%$ des fréquences expérimentales:

  $$
  [p-2 \sigma ; p+2 \sigma] \simeq\left[\frac{1}{6}-2 \times 0,0264 ; \frac{1}{6}+2 \times 0,0264\right] \simeq[0,114 ; 0,219]
  $$

  Autrement dit, avec une marge d'erreur de $5 \%$, l'intervalle $[0,114 ; 0,219]$ contient la fréquence de 6 obtenus.
\end{example}

\newpage
\section{Variables aléatoires}
\subsection*{Introduction}
Une association étudiante souhaite organiser un jeu de dés afin de récolter de l'argent. Les règles du jeu sont les suivantes. Le joueur lance un dé équilibré à six faces.

\begin{itemize}
  \item Il gagne s'il obtient un 6 et l'association lui donne 12 euros.
  \item Il perd dans tous les autres cas et il donne 6 euros à l'association.
\end{itemize}

L'association organise 500 parties de ce jeu. Combien peut-elle espérer gagner en moyenne?

\subsection{Notion de variable aléatoire}
\begin{definition}{}{}
  Une \textbf{variable aléatoire} sur $\Omega$ est une fonction définie sur $\Omega$ à valeur dans $\mathbb{R}$.\\
  À chaque issue de $\Omega$, on associe un nombre réel.
\end{definition}

\begin{definition}{}{}
  Soit $x$ un réel. L'événement « \hide{$\mathrm{X}=x »$} est l'ensemble des issues de $\Omega$ auxquelles on associe le réel $x$.
\end{definition}

\begin{example}
  Avec l'exemple d'introduction :\\
  L'univers est $\Omega=\{1 ; 2 ; 3 ; 4 ; 5 ; 6\}$.\\
  On note $X$ la variable aléatoire correspondant au gain en euros de l'association sur une partie.\\
  L'événement « $\mathrm{X}=6 »$ correspond aux issues 1, 2, 3, 4, 5.\\
  L'événement « $\mathrm{X}=-12$ » correspond à l'issue 6.
\end{example}

\subsection{Loi de probabilité d'une variable aléatoire}
\begin{definition}{}{}
  Pour définir la loi de probabilité d'une variable aléatoire X :

  \begin{itemize}
    \item On liste les valeurs $x_{i}$ possibles pour X ;
    \item On détermine les probabilités $\mathrm{P}\left(\mathrm{X}=x_{i}\right)$.
  \end{itemize}
\end{definition}

\begin{method}{Déterminer la loi de probabilité d'une variable aléatoire}{}
On forme le tableau suivant :

\begin{center}
  \begin{tabular}{|c|c|c|c|c|}
    \hline
    $x_{i}$ & $x_{1}$ & $x_{2}$ & $\ldots$ & $x_{n}$ \\
    \hline
    $P\left(\mathrm{X}=x_{i}\right)$ & $p_{1}$ & $p_{2}$ & $\ldots$ & $p_{n}$ \\
    \hline
  \end{tabular}
\end{center}
\end{method}

\begin{remark}
  Dans le tableau, on a :

  $$
  p_{1}+p_{2}+\ldots+p_{n}=1
  $$
\end{remark}

\begin{example}
  Avec l'exemple d'introduction : Si $X$ est le gain de l'association sur une partie, la loi de probabilité de $X$ est donnée par le tableau suivant :

  \begin{center}
  \renewcommand{\arraystretch}{2}
    \begin{tabular}{|c|c|c|}
      \hline
      $x_{i}$ & -12 & 6 \\
      \hline
      $P\left(\mathrm{X}=x_{i}\right)$ & $\dfrac{1}{6}$ & $\dfrac{5}{6}$ \\
      \hline
    \end{tabular}
  \end{center}
\end{example}

\subsection{Espérance}
\begin{definition}{}{}
  Soit $X$ une variable aléatoire. On appelle \textbf{espérance} de $X$ le nombre :

  $$
  E(X)=p_{1} x_{1}+p_{2} x_{2}+\ldots+p_{n} x_{n}=\hideM{\sum_{i=1}^{n} p_{i} x_{i}}
  $$
\end{definition}

\begin{remark}
  Le mot «espérance» vient du langage des jeux : lorsque $X$ désigne le gain, $E(X)$ est le gain moyen que peut espérer un joueur lorsque le nombre de parties est important.
\end{remark}

\begin{example}
  Avec l'exemple d'introduction :

  $$
  E(X)=p_{1} x_{1}+p_{2} x_{2}=(-12) \times \frac{1}{6}+6 \times \frac{5}{6}=3
  $$

  Cela signifie que le jeu est défavorable au joueur et qu'en moyenne, l'association gagnera 3 euros par partie. Par conséquent, sur 500 parties, elle peut espérer gagner environ 1500 euros.
\end{example}

\subsection{Variance et Écart-type}
\begin{definition}{}{}
  Soit $X$ une variable aléatoire.

  \begin{itemize}
    \item On appelle \textbf{variance} de $X$ le nombre :

  $$
  V(X)=p_{1}\left(x_{1}-E(X)\right)^{2}+p_{2}\left(x_{2}-E(X)\right)^{2}+\ldots+p_{n}\left(x_{n}-E(X)\right)^{2}=\hideM{\sum_{i=1}^{n} p_{i}\left(x_{i}-E(X)\right)^{2}}
  $$

    \item On appelle \textbf{écart-type} de $X$ le nombre $\sigma(X)=\hideM{\sqrt{V(X)}}$.
  \end{itemize}
\end{definition}

\newpage

\begin{example}
  Avec l'exemple d'introduction :

  $$
  \begin{aligned}
    V(X) & =p_{1}\left(x_{1}-E(X)\right)^{2}+p_{2}\left(x_{2}-E(X)\right)^{2} \\
    & =\frac{1}{6}(-12-3)^{2}+\frac{5}{6}(6-3)^{2} \\
    & =45
  \end{aligned}
  $$

  On a ainsi, $\sigma(X)=\sqrt{V(X)}=\sqrt{45} \simeq 6,71$.
\end{example}

\subsection{Transformation affine d'une variable aléatoire}
\begin{proposition}{}{}
  Si $X$ est une variable aléatoire prenant les valeurs $x_{1}, x_{2}, \ldots, x_{n}$ avec les probabilités $p_{1}, p_{2}, \ldots, p_{n}$ et si $a$ et $b$ des réels alors la loi de la variable aléatoire $Y=a X+b$ est donnée par le tableau suivant :

  \begin{center}
    \begin{tabular}{|c|c|c|c|c|}
      \hline
      $y_{i}$ & $a x_{1}+b$ & $a x_{2}+b$ & $\ldots$ & $a x_{n}+b$ \\
      \hline
      $P\left(Y=y_{i}\right)$ & $p_{1}$ & $p_{2}$ & $\ldots$ & $p_{n}$ \\
      \hline
    \end{tabular}
  \end{center}
\end{proposition}

\begin{example}
  Avec l'exemple d'introduction, supposons que l'on multiplie les gains et les pertes par 2 pour chaque partie, on obtient ainsi une nouvelle variable aléatoire $Y=2 X$.
\end{example}

\begin{proposition}{}{}
  Si $X$ est une variable aléatoire et $a$ et $b$ des réels alors

  $$
  E(a X+b)=\hideM{a E(X)+b}
  $$
\end{proposition}

\begin{demo}
  $$
  \begin{aligned}
    E(a X+b) & =p_{1}\left(a x_{1}+b\right)+p_{2}\left(a x_{2}+b\right)+\ldots+p_{n}\left(a x_{n}+b\right) \\
    & =p_{1} a x_{1}+p_{1} b+p_{2} a x_{2}+p_{2} b+\ldots+p_{n} a x_{n}+p_{n} b \\
    & =a\left(p_{1} x_{1}+p_{2} x_{2}+\ldots+p_{n} x_{n}\right)+\left(p_{1}+p_{2}+\ldots+p_{n}\right) b \\
    & =a E(X)+b \quad\left(\operatorname{car} p_{1}+p_{2}+\ldots+p_{n}=1\right)
  \end{aligned}
  $$
\end{demo}

\begin{proposition}{}{}
  Si $X$ est une variable aléatoire et $a$ et $b$ des réels alors

  $$
  V(a X+b)=\hideM{a^{2} V(X)}
  $$
\end{proposition}
\newpage
\begin{demo}
  $V(a X+b)$\\
  $=p_{1}\left(a x_{1}+b-E(a X+b)\right)^{2}+p_{2}\left(a x_{2}+b-E(a X+b)\right)^{2}+\ldots+p_{n}\left(a x_{n}+b-E(a X+b)\right)^{2}$\\
  $=p_{1}\left(a x_{1}+b-a E(X)-b\right)^{2}+p_{2}\left(a x_{2}+b-a E(X)-b\right)^{2}+\ldots+p_{n}\left(a x_{n}+b-a E(X)-b\right)^{2}$\\
  $=p_{1}\left(a x_{1}-a E(X)\right)^{2}+p_{2}\left(a x_{2}-a E(X)\right)^{2}+\ldots+p_{n}\left(a x_{n}-a E(X)\right)^{2}$\\
  $=p_{1}\left(a\left(x_{1}-E(X)\right)\right)^{2}+p_{2}\left(a\left(x_{2}-E(X)\right)\right)^{2}+\ldots+p_{n}\left(a\left(x_{n}-E(X)\right)\right)^{2}$\\
  $=a^{2}\left(p_{1}\left(x_{1}-E(X)\right)^{2}+p_{2}\left(x_{2}-E(X)\right)^{2}+\ldots+p_{n}\left(x_{n}-E(X)\right)^{2}\right)$\\
  $=a^{2} V(X)$
\end{demo}

\begin{proposition}{}{}
  Si $X$ est une variable aléatoire et $a$ et $b$ des réels alors

  $$
  \sigma(a X+b)=|a| \sigma(X)
  $$
\end{proposition}

\begin{demo}
  $$
  \begin{aligned}
    \sigma(a X+b) & =\sqrt{V(a X+b)} \\
    & =\sqrt{a^{2} V(X)} \\
    & =\sqrt{a^{2}} \sqrt{V(X)} \\
    & =|a| \sigma(X)
  \end{aligned}
  $$
\end{demo}

\begin{savoir}{}{}
  \begin{itemize}
    \item Modéliser une situation à l'aide d'une variable aléatoire.
    \item Déterminer la loi de probabilité d'une variable aléatoire.
    \item Calculer l'espérance, la variance et l'écart-type d'une variable aléatoire.
    \item Interpréter l'espérance d'une variable aléatoire.
  \end{itemize}
\end{savoir}

\end{document}